\begin{abstract}
Machine-learning intrusion detectors report near-perfect accuracy on public benchmarks, yet their behaviour in deployment is poorly understood. Deployment brings three challenges at once: traffic from a different network, attack families absent from training, and an analyst team that can review only a fixed number of alerts. A recent review found that no study published between 2024 and 2026 evaluated these concerns together. We present a joint evaluation of four model families on five public benchmarks (CIC-IDS2017, CSE-CIC-IDS2018, ACI-IoT-2023, UNSW-NB15 and ToN-IoT), covering six directed cross-dataset pairs over a harmonised 77-feature flow representation and 47 leave-one-family-out experiments, with discrimination, calibration, base-rate precision and review-budget metrics reported side by side. We also propose a budgeted conformal triage that tests each flow against the hypothesis ``benign and typical of training traffic'' by fusing detector and novelty p-values calibrated on benign flows. In-domain, all models reach AUROC 0.96--0.99, with calibration error at most 0.002 after isotonic recalibration. Across datasets, AUROC falls to 0.50--0.62 and calibration error rises to 0.24--0.45, and recalibration on source data does not repair it. Transferred thresholds either raise the false-positive rate 5--32-fold or detect only 2--3\% of attacks. Within a dataset, 14 of 47 unseen families are detected less than half of the time, 11 of them with over half of their flows given an attack probability below 0.05. Enforcing the review budget with a quantile of unlabeled target traffic keeps the median queue size within 5\% of plan under every shift, whereas transferred thresholds miss it by factors of 0.4 to 54. Fisher fusion of novelty and detector evidence roughly doubles recall at a 1\% budget on the families the detector misses (0.07 to 0.15) at a cost of up to 6 points on families it already detects. Code, data provenance and all results are released.
\end{abstract}
